\section{Модел на заплахите и метод}
\label{sec:threat-model}

Теоретичната основа от глава~\ref{sec:theory} показа при какви предпоставки
конструирането по време на компилация изключва инжекцията. За да бъде това
твърдение проверимо, е необходимо да се определи срещу какъв нападател се
твърди гаранцията, какво точно се защитава и по какви критерии ще се прецени
успехът. Настоящата глава фиксира тези три неща преди представянето на каквито и
да било резултати и описва метода, по който те се проверяват.

\subsection{Защитаван актив}
\label{sec:tm-asset}

Защитаваният актив е \textbf{структурата на командата}, която приложението
изпраща към хранилището за данни. Стойностите в една заявка са данни, а
структурата определя какво хранилището е упълномощено да върне или промени.
Нападател, който превръща {\ttfamily WHERE id = \$1} в
{\ttfamily WHERE id = \$1 OR 1=1}, не е променил данни. Той е променил обхвата на
правомощията на заявката. Всеки клас инжекция от
таблица~\ref{tab:theory-languages} е разновидност на този ход.

\subsection{Модел на нападателя}
\label{sec:tm-attacker}

Нападателят контролира вход, който приложението използва при изграждането на
заявка. Той не контролира изходния код на приложението, компилатора и
драйвера на хранилището. Разграничават се три нива по нарастваща сила.
Разграничението е съществено, защото обичайното обсъждане на инжекциите ги
слива.

\textbf{Ниво 1: нападателят контролира стойности.} Стойност от формуляр, заглавна
част на заявка или съобщение достига заявката. Това е обичайният случай и
единственият, който свързването на параметри покрива.

\textbf{Ниво 2: нападателят контролира идентификатори.} Нападателят избира име на
таблица, на колона или ключ за подреждане. Тук свързването на параметри не
помага, защото нито едно от изследваните хранилища не позволява параметър на
мястото на идентификатор.

\textbf{Ниво 3: нападателят контролира структурата.} Нападателят определя формата
на заявката, например кои условия участват в клаузата \code{WHERE}, чрез
построител на филтри, който сглобява предикати според заявката.

Гаранцията на изследваната библиотека се твърди и за трите нива. На ниво 1 тя
разчита на свързване, а на нива 2 и 3 на това, че съответната възможност
изобщо не е изразима по време на изпълнение (точка~\ref{sec:theory-compile-time}).
Изключение са миграциите, където защитата на ниво 2 е проверка по време на
изпълнение (точка~\ref{sec:enf-ddl}).

\subsection{Извън обхвата}
\label{sec:tm-scope}

Следните заплахи съзнателно са извън твърдяната гаранция и се обсъждат като
ограничения, а не като пропуски:

\begin{itemize}
    \item \textbf{инжекция от втори ред в приложението}: стойност, прочетена от
          хранилището и използвана от приложението като идентификатор другаде, е
          извън видимостта на библиотеката. Собственият случай на библиотеката,
          миграция, която чете имена обратно от хранилището, е в обхвата
          (точка~\ref{sec:enf-ddl});
    \item \textbf{логически грешки в приложението}: липсваща проверка на
          правомощия или прекалено широк филтър не са инжекция;
    \item \textbf{уязвимости в драйверите} (\code{libpq}, \code{hiredis},
          \code{bsoncxx} и останалите), които се приемат за коректни;
    \item \textbf{отказ от услуга} чрез скъпи, но легитимни заявки, който е
          въпрос на капацитет, а не на инжекция.
\end{itemize}

\subsection{Изследователски въпрос и хипотеза}
\label{sec:tm-question}

Работата отговаря на въпроса:

\begin{quote}
\textit{Може ли една абстракция, конструираща заявките по време на компилация, да
елиминира цяло семейство от класове инжекции при хетерогенни хранилища за
данни, и къде точно спира тази гаранция?}
\end{quote}

Хипотезата е, че при изпълнени предпоставки от
точка~\ref{sec:theory-compile-time} никой полезен товар, подаден през
документирания интерфейс, няма наблюдаем ефект върху хранилището, докато същият
полезен товар, подаден в наивно сглобена заявка, има. Втората част на въпроса,
къде спира гаранцията, е равностойна на първата. Работа, която отчита само
успехи, не позволява да се прецени валидността на твърденията си.

\subsection{Критерии за успех}
\label{sec:tm-criteria}

Критериите са определени преди измерванията и не са променяни след тях.

\begin{table}[H]
\centering
\caption{Критерии за успех}
\label{tab:tm-criteria}
\begin{tabularx}{\textwidth}{L{1.0cm}Y}
\toprule
\textbf{№} & \textbf{Критерий} \\
\midrule
К1 & Никой полезен товар от корпуса, подаден през документирания интерфейс, не
     предизвиква наблюдаем ефект: върнати редове, които не би трябвало да
     съществуват, промяна в данните или схемата, или принудително изпълнение. \\
К2 & Всеки опит идентификатор или структура на заявката да зависи от стойност по
     време на изпълнение се отхвърля от компилатора, и то по причината, която се
     проверява, а не по случайна друга. \\
К3 & Наивната заявка, сглобена от същите полезни товари, е компрометирана поне
     веднъж при същата постановка. Без това изпълнението на К1 не доказва нищо. \\
К4 & Границите на гаранцията са назовани и, където е възможно, демонстрирани. \\
\bottomrule
\end{tabularx}
\end{table}

Критерий К3 заслужава коментар. Сравнение, при което нито защитеният, нито
незащитеният път са засегнати, не различава защитата от неефективен корпус.
Наивната заявка е съзнателно уязвима и служи за база, спрямо която резултатът на
библиотеката придобива смисъл. Ако тя не бъде компрометирана, тестът се отчита
като неуспешен.

\subsection{Метод на оценката}
\label{sec:tm-method}

Методът съчетава три независими източника на доказателства, всеки от които
покрива различна предпоставка.

\textbf{Отрицателни компилационни тестове} (за К2). Всеки тест е програма, която
се опитва да направи идентификатор или структура зависими от стойност по време на
изпълнение и трябва да бъде отхвърлена от компилатора. Такъв тест сам по себе си
е слаб: той би преминал и при грешка, несвързана с проверяваното свойство,
например липсващ заглавен файл. Затова се прилагат две мерки. Първо, всяко
отхвърляне трябва да съответства на образец за очакваната диагностика. Второ,
наборът съдържа \textbf{контролни случаи}, които използват същата постановка и
трябва да се компилират. Ако постановката е повредена, контролните случаи също
се провалят.

\textbf{Диференциално тестване} (за К1 и К3). Всеки приложим полезен товар се
подава по два пътя, които задават един и същ въпрос към една и съща база данни:
през наивно сглобена заявка и през библиотеката. За всяко изпълнение се използва
нова база данни. Резултатът се определя от \textbf{оракул}, който наблюдава само
ефекти: върнати редове, промяна на състоянието и принудително прекратяване на
изпълнението.

\textbf{Анализ на изходния код} (за предпоставка 2 от
точка~\ref{sec:theory-compile-time} и за К4). За всеки от седемте конектора се
установява как точно стойността достига драйвера.

\subsection{Какво методът не доказва}
\label{sec:tm-not-proven}

Методът има граници, които се посочват предварително, за да не се надценяват
резултатите.

\begin{itemize}
    \item Оракулът вижда само ефекти. Сляпа атака, която променя логиката на
          заявката, но връща нула редове, не се отчита като компрометиране.
          Следователно броят на компрометиранията на наивния път е \textbf{долна
          граница}. Това изместване работи срещу твърдението на работата, а не в
          негова полза.
    \item Диференциалният тест използва една форма на заявка за хранилище. Друга форма би
          могла да позволи други атаки срещу наивния път.
    \item Корпусът е краен. Отсъствие на успешна атака срещу библиотеката не е
          доказателство за отсъствие на всякаква атака. Доказателството за това се
          търси в аргумента от точка~\ref{sec:theory-compile-time} и в
          компилационните тестове, а не в корпуса.
    \item Компилационните тестове се изпълняват с GCC и Clang. Поведението на
          други компилатори не е проверено.
\end{itemize}
